\section{Empty original entries in the entire recursive canonical value}
The ordered boundary is part of the actual marked ideal. Inserting an empty
original entry changes its indices, even though it changes neither its ideal
nor its cosupport. The comparison below retains this change throughout the
entire recursive value. It uses the actual masks and histories of the masked
Step-One construction in W\l odarczyk's marked-ideal algorithm. The explicit
recursive index map describes how the retained indexed boundary adapts the
insertion of isomorphisms in the original resolution history.

\begin{definition}[The recursive empty-entry value map]
\label{hir:original-invariant-empty-map}
\source{wlodarczyk05-hironaka,bierstone-milman08-functoriality}
\lean{Hironaka.HironakaOriginalInvariantEmptyInsertion.shiftFamily}
\leanok
\uses{hir:canonical-centres,bcalc:monomial-invariant-empty-member}
Let $F_{0,j}$ be the identity on the zero-dimensional value type. For the
successor value type, $F_{k+1,j}$ retains its first, nonmonomial-order entry.
On its monomial branch it retains bottom and applies the increasing index
insertion map at $j$ to the finite index component of the monomial value. On its positive-order
branch it retains infinity; on a finite Step-One value it retains the original
level and applies $F_{k,0}$ to the recursive lower value, retaining infinity
there as well. Thus all recursive monomial index components after the first
one use the insertion map at zero.
\end{definition}

\begin{lemma}[The recursive map is a strict order embedding]
\label{hir:original-invariant-empty-map-strict}
\source{wlodarczyk05-hironaka,bierstone-milman08-functoriality}
\lean{Hironaka.HironakaOriginalInvariantEmptyInsertion.strictMono_shiftFamily,Hironaka.HironakaOriginalInvariantEmptyInsertion.shiftFamilyEmbedding}
\leanok
\uses{hir:original-invariant-empty-map,bcalc:monomial-invariant-empty-member}
Each $F_{k,j}$ is strictly increasing and consequently reflects strict order,
weak order, and equality on the entire recursive value type.
\end{lemma}
\begin{proof}
Induct on $k$. The identity has the assertion at zero. At a successor, the
monomial index insertion preserves and reflects the colexicographic order
of finite sets, and leaves the other monomial entries unchanged. Extending a
strictly increasing map by a fixed bottom or top preserves strict order.
The induction therefore applies to the lower-value map, its top extension,
and the Step-One lexicographic pair with unchanged level. Extend that map
again by a fixed top. The sum order puts every monomial branch value below
every positive branch value, so applying these two increasing branch maps
preserves strict order on the sum. Finally the outer lexicographic pair has
an unchanged first entry. If two first entries differ their comparison is
unchanged; if they agree their second-entry comparison is preserved. This
proves the successor assertion. A strictly increasing map between linear
orders reflects both orders and is injective, giving all three reflections.
\end{proof}

\begin{theorem}[The entire actual invariant commutes with empty insertion]
\label{hir:original-invariant-empty-actual}
\source{wlodarczyk05-hironaka,bierstone-milman08-functoriality}
\lean{Hironaka.HironakaOriginalInvariantEmptyInsertion.invFamily_insertEmpty}
\leanok
\uses{hir:original-invariant-empty-map,hir:positive-member-empty-actual,hir:positive-member-empty-bound,hir:member-future-boundary,hir:step-one-empty-insertion-masked-zero,hir:step-one-empty-insertion-run,hir:original-positive-companion-empty-run}
For every $k$, actual marked ideal $m$, valid original insertion position $j$,
genuine blowup tower $T$, and endpoint $y$, the actual canonical value satisfies
\[
 \operatorname{inv}_{k}(\operatorname{insertEmpty}_{j}m,T,y)
 =F_{k,j}(\operatorname{inv}_{k}(m,T,y)).
\]
\end{theorem}
\begin{proof}
Induct on $k$. At zero the invariant is constant. At a successor the actual
nonmonomial orders agree at every stage. If the endpoint has order zero, the
whole transformed boundary has the same insertion at $j$, so its monomial
value is precisely the monomial branch of $F_{k+1,j}$.
For positive order, the actual order regimes agree along every stage, hence
their computed runs and genuine head-tail splits agree. The actual order
masks agree and their actual companion marked ideals have the same insertion
at $j$. The true Step-One original levels, specialness, regimes, and runs
also agree. Special endpoints give identical infinite Step-One values.
Otherwise the actual lower masks agree. At level zero, the masked contact
record has the same insertion at $j$; its entire original prefix has empty
traces on this actual mask because the zero mask requires avoidance of every
actual original member. Pulling each of these members back to the mask
therefore gives the empty set. At positive level the genuine member restriction
computes the new insertion index $c$ as the number of retained entries before
$j$. Its whole retained original prefix is empty on the actual mask, and
$c$ lies inside that prefix. Apply the induction to the real lower record
and genuine restricted tail. At either level the positions selected in the current monomial branch
lie beyond the empty prefix; insertion at $j$, respectively $c$, therefore
acts on those positions exactly as insertion at zero. The
positive branch of $F_{k,c}$ is already independent of $c$. Thus the actual
recursive lower value is $F_{k,0}$ of its old value. Extending through both
actual infinity cases gives the stated positive branch of $F_{k+1,j}$.
All masks and runs here are those computed by the original value; no
comparison, selected-index, or history alignment is assumed.
\end{proof}

\begin{corollary}[All actual value comparisons and maximum loci are retained]
\label{hir:original-invariant-empty-comparisons}
\source{wlodarczyk05-hironaka,bierstone-milman08-functoriality}
\lean{Hironaka.HironakaOriginalInvariantEmptyInsertion.invFamily_insertEmpty_comparisons,Hironaka.HironakaOriginalInvariantEmptyInsertion.maxLocus_invFamily_insertEmpty}
\leanok
\uses{hir:original-invariant-empty-actual,hir:original-invariant-empty-map-strict}
At fixed $k$ and insertion position $j$, all comparisons of actual values
before and after insertion agree, including comparisons between different
genuine histories. In particular every maximum locus on any set of endpoints
of a fixed history agrees.
\end{corollary}
\begin{proof}
Rewrite each actual inserted value by the preceding theorem and apply the
three order reflections of the same strictly increasing map $F_{k,j}$.
A point belongs to a maximum locus exactly when it belongs to its comparison
set and dominates every value on that set. Weak-order reflection identifies
every one of these comparisons, proving equality of the whole maximum loci.
\end{proof}

\begin{corollary}[Entire actual canonical histories are retained]
\label{hir:original-invariant-empty-canonical-histories}
\source{wlodarczyk05-hironaka,bierstone-milman08-functoriality}
\lean{Hironaka.HironakaOriginalInvariantEmptyInsertion.isCanonical_invFamily_insertEmpty_iff,Hironaka.HironakaOriginalInvariantEmptyInsertion.isCanonicalSegment_invFamily_insertEmpty_iff}
\leanok
\uses{hir:original-invariant-empty-comparisons,bcalc:monomial-invariant-empty-member}
For each actual marked ideal and valid empty original insertion, a genuine
history is canonical for the inserted datum if and only if it is canonical
for the old datum. This holds both when empty centres are permitted and when
every centre must be nonempty.
\end{corollary}
\begin{proof}
Induct on the genuine history. At length zero both canonical predicates hold.
At a successor, their preceding histories agree by induction. The whole
transformed cosupports agree because insertion retains the carrier, ideal,
and marking at every stage. The preceding maximum-locus corollary therefore
identifies the actual maximum centres at the preceding stage. The successor
centre is the same set for both inputs, so its emptiness or nonemptiness is
also the same. Thus every clause of either recursive canonical-history
definition agrees, giving both equivalences.
\end{proof}
